\documentclass{beamer}
\usepackage{amsmath}
\usepackage{braket}
\usepackage{color}
\usepackage{bbold}
\usepackage{graphicx} % Required for inserting images
\usepackage{ amssymb } % for check marks
\usepackage{svg}
\usepackage{tikz}

\usetikzlibrary{quantikz2}

\usetheme{Antibes}

\title{Steane Type Error Correction}
\subtitle{Long Update}
% \author{Leonard Lux}
\date{09.06.2026}

\begin{document}

\tikzset{ noisy/.style={fill=red!20}
}

% color code
\tikzset{ red/.style={fill=red!20}}
\tikzset{ green/.style={fill=green!20}}
\tikzset{ blue/.style={fill=blue!20}}

% Title page
\begin{frame}
\titlepage
% \begin{center}
% Supervised by:\newline 
% \end{center}
% \begin{center}
% Luis Colmenarez, Maarten Wegewijs \& Markus Müller
% \end{center}
\end{frame}



\section{Surface Code}
\begin{frame}{Encoding: \textbf{Surface Code}}
    \begin{columns}
    \begin{column}{0.5\textwidth}
        \begin{itemize}
            \setlength{\itemsep}{8pt}
            \item Stabilizer code
            \item Transversal $\overline{\mathrm{CNOT}}$ 
        \end{itemize}
    \end{column}
    \begin{column}{0.5\textwidth}  
        \begin{center}
        \includegraphics[width=0.9\textwidth]{sketch_surface_code.png}
        \end{center}
    \end{column}
    \end{columns}
\end{frame}


\section{Steane Type Error Correction}
\begin{frame}{Steane Type Error Correction}
    \begin{columns}
    \begin{column}{0.6\textwidth}
    \begin{block}{The Circuit}
        % Simple Steane Type EC
        \begin{quantikz}
            \lstick{$\ket{\Psi}_L$}&\qwbundle{n}&\targ{}&&\ctrl{2}&&\\%&&\gate{Z_R}&\gate{X_R}&\\
            \lstick{$\ket{0}_L$}&\qwbundle{n}&\ctrl{-1}&\meter{X}&\setwiretype{c}&&\\%\gate{Decoder}\wire[u][1]{c}\\
            \lstick{$\ket{+}_L$}&\qwbundle{n}&&&\targ{}&\meter{Z}&\setwiretype{c}%&\gate{Decoder}\wire[u][2]{c}
        \end{quantikz}
    \end{block} 
    \end{column}
    \begin{column}{0.4\textwidth}
    \begin{itemize}
        \item Fault-tolerant  
        \item Single-shot 
        \item Not symmetric 
    \end{itemize} 
    \end{column}    
    \end{columns}
\end{frame}


% Slide  (Error Models)

\section{Error Model}

\subsection{Detector Error Model} 
\begin{frame}{Circuit Level Noise}

    \begin{quantikz}
        \lstick{$\ket{\Psi}$}&\targ{}&\ctrl{2}&&\\
        \lstick{$\ket{0}$}&\ctrl{-1}&&\meter{$X$}&\setwiretype{c}\\
        \lstick{$\ket{+}$}&&\targ{}&\meter{$Z$}&\setwiretype{c}\\
    \end{quantikz}

    \begin{quantikz}
        % Data qubit
        \lstick{$\ket{\Psi}$}&\gate[1,style={noisy},]{D}&   % 1: init
        \targ{}&\gate[1,style={noisy}]{D2}\wire[d][1]{q}&      % 2: 1. CNOT 
        \ctrl{2}&\gate[1,style={noisy}]{D2}\wire[d][2]{q}&     % 3: 2. CNOT
        % &&                                                      % 4: 
        &&\\                                                    % 5: 
        % 0 ancilla qubit
        \lstick{$\ket{0}$}&\gate[1,style={noisy},]{D}&      % 1: init
        \ctrl{-1}&\gate[1,style={noisy}]{D2}&                  % 2: 1. CNOT
        &&                                                      % 3: 2. CNOT
        % \gate[]{H}&                                             % 4: 1. change basis aka H
        % & 
        % \gate[1,style={noisy},]{D_1}&                           %    2: Error (No errors on H)
        \gate[1,style={noisy},]{Z}&\meter{$X$}&\setwiretype{c}\\% 5: stabilizer measurement
        % p ancilla qubit 
        \lstick{$\ket{+}$}&\gate[1,style={noisy},]{D}&      % 1: init
        &&                                                      % 2: 1. CNOT
        \targ{-2}&\gate[1,style={noisy}]{D2}&                  % 3: 2. CNOT
        % &&                                                      % 4: 
        \gate[1,style={noisy},]{X}&\meter{$Z$}&\setwiretype{c}\\% 5: stabilizer measurement
    \end{quantikz}

\end{frame}

\begin{frame}{1. Last Time - Syndrome Channel}

    \begin{quantikz}
        % Data qubit
        \lstick{$\ket{\Psi}$}&\gate[1,style={green},]{D}&   % 1: init
        \targ{}&\gate[1,style={green}]{D2}\wire[d][1]{q}&      % 2: 1. CNOT 
        \ctrl{2}&\gate[1,style={blue}]{D2}\wire[d][2]{q}&     % 3: 2. CNOT
        % &&                                                      % 4: 
        &&\\                                                    % 5: 
        % 0 ancilla qubit
        \lstick{$\ket{0}$}&\gate[1,style={green},]{D}&      % 1: init
        \ctrl{-1}&\gate[1,style={green}]{D2}&                  % 2: 1. CNOT
        &&                                                      % 3: 2. CNOT
        % \gate[]{H}&                                             % 4: 1. change basis aka H
        % & 
        % \gate[1,style={noisy},]{D_1}&                           %    2: Error (No errors on H)
        \gate[1,style={noisy},]{Z}&\meter{$X$}&\setwiretype{c}\\% 5: stabilizer measurement
        % p ancilla qubit 
        \lstick{$\ket{+}$}&\gate[1,style={blue},]{D}&      % 1: init
        &&                                                      % 2: 1. CNOT
        \targ{-2}&\gate[1,style={blue}]{D2}&                  % 3: 2. CNOT
        % &&                                                      % 4: 
        \gate[1,style={blue},]{X}&\meter{$Z$}&\setwiretype{c}\\% 5: stabilizer measurement
    \end{quantikz}

    \begin{quantikz}
        % Data qubit
        \lstick{$\ket{\Psi}$}&
        \targ{}&
        \ctrl{2}&
        \meter{trace}
        % &                                                      
        \\                                                    
        % 0 ancilla qubit
        \lstick{$\ket{0}$}&
        \ctrl{-1}&
        &
        &
        \gate[1,style={green}]{D'}\wire[d][1]{q}&
        \gate[1,style={red},]{Z}&\meter{$X$}&\setwiretype{c}\\% 5: stabilizer measurement
        % p ancilla qubit 
        \lstick{$\ket{+}$}&
        &
        \targ{-2}&
        &
        \gate[1,style={green}]{D'}&
        % &&                                                      % 4: 
        \gate[1,style={blue},]{X}&\meter{$Z$}&\setwiretype{c}\\% 5: stabilizer measurement
    \end{quantikz}
\end{frame}


\begin{frame}{Syndrome Channel: Product Channel}

\begin{itemize}
    \item Solutions: cutting a wire(s) or trivial
\end{itemize}
\begin{align*}
    p_D^c &= 0 \ \ \ \rightarrow  \  p_{t,ij} =  p_{x,i}^c \cdot p_{z,j}^c\\
    p_D^c &= 3/4 \rightarrow\   p_{t,ij} = \frac{1}{4} \\
    p_x^c &= 1/2 \rightarrow \  p_{t,ij} =  \frac{1}{3}p_D^c + \frac{1}{2}\left(1-\frac{4}{3}p_D^c\right)p_{z,j}^c\\
    p_z^c &= 1/2 \rightarrow \  p_{t,ij} =  \frac{1}{3}p_D^c + \frac{1}{2}\left(1-\frac{4}{3}p_D^c\right)p_{x,i}^c
\end{align*}

\end{frame}

\begin{frame}{Syndrom Channel: Depolarizing  Channel}
\begin{itemize}
    \item Solutions: cutting both wires or trivial
\end{itemize}
\begin{align*}
    p_D^c &= 3/4 \rightarrow p_{t,ij} = 1/4\\
    p_x^c = p_z^c &= 1/2\rightarrow p_{t,ij} = 1/4\\
    p_x^c = p_z^c &= 0 
\end{align*}

\end{frame}

% Put here coherent stuff
\begin{frame}{Syndrome- and Data Channel: Measurement errors}
    \begin{itemize}
        \item complete channel (up to first order in $p$)
    \end{itemize}

    \begin{equation*}
         \hat{N}(\rho) = \sum_i \sum_j p_{ij} \hat{E}_i^{s} \hat{E}_j^{\ket{\Psi}} + \mathcal{O}(p^2)
    \end{equation*}
    
    \begin{itemize}
        \item no measurement errors:     
    \end{itemize}
    \begin{align*}
        \sum_i p_{ii} &= 1- \frac{2}{3} p_{sp\ket{0}} -  p_{sp\ket{+}} - p_{m\ket{0}}- p_{m\ket{+}} - \frac{8}{15} p_{D2\ket{0}}- \frac{4}{5} p_{D2\ket{+}}  \\
        &= 1- 5 p 
    \end{align*}


\end{frame}



\begin{frame}

    \begin{itemize}
        \item no $X$ measurement errors:     
    \end{itemize}
    \begin{align*}
        \sum_i \sum_{j \in \{1,Z\}} p_{i,ij} &= 1- \frac{2}{3} p_{sp\ket{+}} - p_{m\ket{0}} - \frac{8}{15} p_{D2\ket{0}} \\
        &= 1- \frac{11}{5} p 
    \end{align*}

    \begin{itemize}
        \item no $Z$ measurement errors:     
    \end{itemize}
    \begin{align*}
        \sum_i \sum_{j \in \{1,X\}} p_{i,ij} &= 1- \frac{2}{3} (p_{sp\ket{+}} + p_{sp\ket{0}}) - p_{m\ket{+}} - \frac{8}{15} (p_{D2\ket{0}}+ p_{D2\ket{0}} ) \\
        &= 1- \frac{17}{5} p 
    \end{align*}

\end{frame}




\section{Numerical Precision}
\begin{frame}{2. Numerical Precision Problem of ML}
\begin{figure}
    \centering
    \includegraphics[width=0.6\textwidth]{plots/ml_test_example.pdf}
\end{figure}
\begin{itemize}
    \item Assumption: numerical precision at fault
    \begin{itemize}
        \item higher $d$ $\rightarrow$ smaller coset probability $\rightarrow$ higher instability
    \end{itemize}
    \item Test:
    \begin{itemize}
        \item Increase/Decrease precision  $\rightarrow$ stability ?
        \item $\log$ of coset probability 
    \end{itemize}
\end{itemize}
\end{frame}

\begin{frame}

\begin{figure}
    \centering
    \includegraphics[width=0.6\textwidth]{plots/ml_test_different_approaches.pdf}
\end{figure}
\begin{itemize}
    \item reduced precision (32-bit) $\rightarrow$ earlier instability
    \begin{itemize}
        \item 16-bit float and 128-bit float not implemented 
    \end{itemize}
    \item $\log$ does not make a difference 
    \item instability appears deterministic  
\end{itemize}
\end{frame}

\begin{frame}{Code  Capacity}
\begin{columns}

    \begin{column}{0.5\textwidth}
    \begin{figure}
        \centering
        \includegraphics[width=0.8\textwidth]{plots/ml_test_circ_d_3.pdf}
    \end{figure}
    \begin{figure}
        \centering
        \includegraphics[width=0.8\textwidth]{plots/ml_test_circ_d_5.pdf}
    \end{figure}
    \end{column}    

    \begin{column}{0.5\textwidth}
    \begin{figure}
        \centering
        \includegraphics[width=0.8\textwidth]{plots/ml_test_circ_d_9.pdf}
    \end{figure}
    \begin{figure}
        \centering
        \includegraphics[width=0.8\textwidth]{plots/ml_test_circ_d_15.pdf}
    \end{figure}
    \end{column}
\end{columns}
\end{frame}



\section{Repeated Syndrome Extraction}
\begin{frame}{3. Reference - Repeated Syndrome Extraction}
\begin{columns}
    \begin{column}{0.5\textwidth}
        
    \begin{itemize}
        \item 'standard' syndrome readout
        \item repeated $d$-times per round
        \item Decoder:
        \begin{itemize}
            \item MWPM
            \item all rounds + final faultless measurement (FT check) 
        \end{itemize} 
    \end{itemize} 
    \end{column}

    \begin{column}{0.5\textwidth}
        Standard: 
        \begin{quantikz}
            \lstick{$q_1$} & \ctrl{4} & \qw      & \qw      & \qw      & \qw \\
            \lstick{$q_2$} & \qw      & \ctrl{3} & \qw      & \qw      & \qw \\
            \lstick{$q_3$} & \qw      & \qw      & \ctrl{2} & \qw      & \qw \\
            \lstick{$q_4$} & \qw      & \qw      & \qw      & \ctrl{1} & \qw \\
            \lstick{$a$}   & \targ{}  & \targ{}  & \targ{}  & \targ{}  & \meter{}
        \end{quantikz}
    \end{column}

\end{columns}
\end{frame}

\begin{frame}
    \begin{figure}
        \centering
        \includegraphics[width=0.65\textwidth]{plots/multiple_rounds_ref.pdf}
    \end{figure}

    \begin{table}
    \begin{tabular}{lcc}
    \hline
    Type & Obs. & Threshold (\%) \\
    \hline
    Steane (ML)   & Z & $1.988 \pm 0.041$ \\
    Steane (ML)   & X & $1.982 \pm 0.038$ \\
    \hline
    Rep. (MWPM)   & Z & $0.852 \pm 0.067$ \\
    Rep. (MWPM)   & X & $0.897 \pm 0.019$ \\
    \end{tabular}
    \label{tab:final_results}
    \end{table} 
\end{frame}

\begin{frame}
    \begin{center}
        Thanks for listening!
    \end{center}  
\end{frame}



\end{document}